% !TEX root = ../monograph.tex
\chapter{Relativistic reconstruction in every phase}
\label{ch16:reconstruction}

We now pass from separated Euclidean correlations to relativistic fields.
The order of the argument matters. Bounded slab insertions give analytic
functions of complex time and their joint tempered boundary values.
Lorentz covariance and spectral support then supply a full spacetime
tube, in which Euclidean permutation symmetry becomes locality. Only
after these steps do we form the Hilbert space of field words.

Throughout this chapter we use the conclusions of
Chapters~\ref{ch14:reflection} and~\ref{ch15:phases}: the coordinate
reflection spaces, exact Euclidean symmetries and common phase field
calculus. Write $\Xi$ for the compact phase space, $\nu$ for its
full-support probability measure, and $S^\xi$ for its correlations.
The finite-dimensional spaces $F_d$ have compatible inclusions; their
algebraic direct limit is denoted by $F$. All estimates below concern a
fixed finite list of labels. Uniformity means uniformity in $\xi$, not in
the number or degree of those labels.

The underlying reconstruction principles are those of Wightman and
Osterwalder--Schrader \cite{Wightman1956,OsterwalderSchrader1973,OsterwalderSchrader1975}.
Here the available Euclidean data are smooth functions off the diagonals,
with bounded slab operators. We therefore verify the boundary construction
directly; an unproved extension of these functions across Euclidean
collisions is not an input.

\section{Complex time and a uniform boundary distribution}
\label{ch16:boundary-section}

Use coordinate four as time. In the phase reflection space let
$E_\xi(s)=e^{-sH_\xi^{\rm OS}}$ and $U_\xi(y)$ be the heat and spatial
translation operators. The spectral theorem extends $E_\xi$ to
$\operatorname{Re}s>0$ as a holomorphic contraction. Let $B_\varepsilon^\xi(A)$
be the slab operator of width $2\varepsilon$ with the single label $A$
placed at its center. The phase version of the slab bound gives
\begin{equation}
 \|B_\varepsilon^\xi(A)\|\le C_A\varepsilon^{-q_A},\qquad
 B_{\varepsilon+u}^\xi(A)
       =E_\xi(u)B_\varepsilon^\xi(A)E_\xi(u),
 \quad 0<\varepsilon\le1,
 \label{ch16:slabs}
\end{equation}
where $u\ge0$ in the padding identity. Linearity and compatibility under
$F_d\hookrightarrow F_{d+1}$ follow by testing against all exterior words.

For $z_j=t_j+iy_j$ with $y_1<\cdots<y_n$, set
\[
 G_\varepsilon^\xi(z,x)
       =E_\xi(-iz-2\varepsilon)U_\xi(x).
\]
For $4\varepsilon<\min_j(y_{j+1}-y_j)$ define
\begin{equation}
 \mathsf F_n^\xi(\boldsymbol A;z,\boldsymbol x)
  =\left\langle\Omega_\xi^{\rm OS},
    B_\varepsilon^\xi(A_1)
    \prod_{j=2}^n\left[
       G_\varepsilon^\xi(z_j-z_{j-1},x_j-x_{j-1})
       B_\varepsilon^\xi(A_j)\right]\Omega_\xi^{\rm OS}\right\rangle.
 \label{ch16:heat-function}
\end{equation}
The product is ordered from left to right with increasing $j$.
Padding and the semigroup law show that this expression is independent
of $\varepsilon$: replacing it by $\varepsilon+u$ moves two factors
$E_\xi(u)$ onto each intervening gap, while the exterior factors fix the
vacuum. At purely imaginary times it is $S^\xi$ at the corresponding
Euclidean points. It is holomorphic in the time variables and continuous
in the real spatial variables. No norm continuity of the spatial
unitaries is required.

Contraction and \eqref{ch16:slabs}, with $\varepsilon=\delta/8$, give
\begin{equation}
 |\mathsf F_n^\xi(t+iy,x)|\le C\delta^{-m},\qquad
 0<\delta\le1,\quad y_{j+1}-y_j\ge\delta,
 \label{ch16:gap-bound}
\end{equation}
uniformly in all real positions and in $\xi$.
This is the quantitative regularity needed for reconstruction.

\begin{proposition}[Joint tempered boundary]
\label{ch16:boundary}
For every label tuple there is a tempered distribution
$W_n^\xi(\boldsymbol A)$ in all $4n$ real coordinates, obtained as the
boundary of \eqref{ch16:heat-function} along arbitrary strictly ordered
heights tending to zero. The distributions obey one common Schwartz
seminorm bound, depend continuously on $\xi$ when tested, and have the
mixed correlation distribution as their $\nu$-average. The same assertions
hold at fixed weakly ordered imaginary heights, including tied heights.
\end{proposition}
\begin{proof}
For a joint Schwartz test function $f$ put
$\psi_y^\xi(f)=\int\mathsf F_n^\xi(t+iy,x)f(t,x)\,dt\,dx$.
Let $r=(0,1,\ldots,n-1)$ and $D=-i\sum_jr_j\partial_{t_j}$ acting on test functions.
Cauchy's formula, \eqref{ch16:gap-bound}, and integration by parts give
$\frac d{ds}\psi_{y+sr}^\xi(f)=\psi_{y+sr}^\xi(Df)$.
Taylor's formula, evaluated backwards from $s=1$, is the exact identity
\begin{equation}
 \psi_y^\xi(f)=
 \sum_{k=0}^m\frac{(-1)^k}{k!}\psi_{y+r}^\xi(D^kf)
 -\int_0^1\frac{(-s)^m}{m!}
                   \psi_{y+sr}^\xi(D^{m+1}f)\,ds.
 \label{ch16:taylor-boundary}
\end{equation}
The integrand is bounded by $C\|D^{m+1}f\|_1$ for every weakly ordered
$y$: the factor $s^m$ cancels the worst gap singularity. The right-hand
side therefore defines the extension to all such $y$, continuously by
dominated convergence. Its bound is a fixed finite sum of $L^1$ norms
of derivatives of $f$, hence a Schwartz seminorm. In particular there
is no restriction on ratios of the heights approaching zero.

For strictly ordered heights, phase continuity follows from the normal
family theorem in the time variables: every convergent phase sequence
has locally convergent analytic subsequences, and agreement at imaginary
times identifies each limit with the limiting phase by the identity
theorem. Local boundedness is uniform by \eqref{ch16:gap-bound}.
Dominated convergence after testing gives continuity of $\psi_y^\xi$.
The same uniqueness argument continues the barycenter identity from
imaginary times. Formula~\eqref{ch16:taylor-boundary} transfers both
claims to weakly ordered heights and to the boundary. Arity zero and
one require no gap argument and are covered by their constant bounds.
\end{proof}

\section{Lorentz covariance and the spectral cone}
\label{ch16:covariance-section}

Let $s_4$ be reflection of Euclidean coordinate four. The field adjoint
on $F$ is the antilinear involution
\[
                       A^*=R(s_4)\overline A.
\]
If $L_{ab}$ is the infinitesimal Euclidean rotation action on labels,
the spatial rotation generators are $L_{ij}$ and the boost generators
are $-iL_{4i}$. Their commutation relations give the proper orthochronous
Lorentz action $\rho$. It descends from the covering group because a full
spatial turn already acts trivially under $R$. The above involution
commutes with $\rho$: the sign from conjugating $i$ cancels the reflection
sign on $L_{4i}$. We use spacetime coordinates $(x^1,x^2,x^3,t)$,
metric $\eta=\operatorname{diag}(-1,-1,-1,1)$, and
\[
 \overline V_+=\{(p,E):E\ge|p|\}.
\]
Fourier duality uses the ordinary coordinate pairing $p\cdot x+Et$.
The dual Lorentz action is the inverse transpose, which also preserves
$\overline V_+$.

Differentiate Euclidean rotational covariance at imaginary times and
continue the identity using the Cauchy--Riemann equations. The result
is the infinitesimal Lorentz covariance identity for $\mathsf F_n^\xi$.
Testing and taking its joint boundary proves the distributional identity:
terms containing a height tend to zero because the tested boundary
families have the common seminorm bound of Proposition~\ref{ch16:boundary}.
Integrating the finite-dimensional label action and the corresponding
flow on Schwartz test functions proves Lorentz covariance of $W_n^\xi$.
Translations and spatial rotations continue directly.

Use the Fourier convention
\[
 \widehat f(X)=\int e^{-iX\cdot p}f(p)\,dp,
 \qquad \widehat W(f)=W(\widehat f).
\]
For a compactly supported test function $b$ in momentum space, differentiating simultaneously
the imaginary height and the factor multiplying $b$ gives
\begin{equation}
 \widehat W_n^\xi(b)
   =\psi_y^\xi\!\left(\widehat{e^{\sum_jy_jp_j^4}b}\right).
 \label{ch16:laplace}
\end{equation}
Indeed the two differentiated terms cancel, and their common value at
height zero is the left side. Compact frequency support makes every
intermediate test function a Schwartz function.

Translation invariance restricts support to $\sum_jp_j=0$. Put
$q_j=\sum_{k=j}^np_k$. If the support of $b$ lies in $q_j^4<0$,
let the corresponding imaginary gap tend to infinity in
\eqref{ch16:laplace}, keeping the other gaps positive. Exponential decay
of the weighted test function and the uniform bound at gaps at least one show that its
pairing vanishes. A partition of unity and Lorentz covariance now give
\begin{equation}
 \operatorname{supp}\widehat W_n^\xi\subset K_n,
 \qquad K_n=\left\{p:\sum_jp_j=0,\quad q_j\in\overline V_+ (2\le j\le n)\right\}.
 \label{ch16:cone-support}
\end{equation}
To justify the last step, a vector outside $\overline V_+$ can be made
to have negative time component by a proper orthochronous Lorentz
transformation. Compact subsets of the complement admit finitely many
such neighborhoods.

\section{The spacetime tube and locality}
\label{ch16:locality-section}

The boundary distribution just constructed is denoted by $W_n^\xi$;
$K_n$ in \eqref{ch16:cone-support} is its momentum support cone. Define
\[
 C_n=\{Y:Y_{j+1}-Y_j\in\overline V_+\},\qquad
 \mathcal T_n=\{X+iY:Y_{j+1}-Y_j\in V_+\}.
\]
For $p\in K_n$, telescoping gives
$p\cdot Y=\sum_{j=2}^nq_j\cdot(Y_j-Y_{j-1})\ge0$.
For strictly interior $Y$ this is bounded below by $c_Y|p|$, locally
uniformly in $Y$. One way to see the strict margin is to choose a ball
of radius $2c_Y$ around $Y$ inside the strict-gap cone and apply positivity
to $Y-c_Yp/|p|$.

Choose a smooth cutoff $\chi$ equal to one on the distance-one neighborhood
of $K_n$, supported in its distance-three neighborhood, with all derivatives
bounded. Convolution of the distance-two neighborhood's indicator with
a bump of radius $1/2$ is one such choice. Then
\begin{equation}
 \mathcal W_n^\xi(\boldsymbol A;Z)
   =(2\pi)^{-4n}\widehat W_n^\xi
                    \left(\chi(p)e^{ip\cdot Z}\right),\qquad Z\in\mathcal T_n,
 \label{ch16:full-tube}
\end{equation}
is well defined: the tested function and all its derivatives decay
exponentially on the support of $\chi$. It is holomorphic into Schwartz
space before the distribution is applied. It agrees with
\eqref{ch16:heat-function} at strictly ordered scalar heights, by
\eqref{ch16:laplace} and compact-frequency approximation.

For $Y\in C_n$, the multiplier $\chi e^{-p\cdot Y}$ and its derivatives
are uniformly bounded on compact sets of $Y$. In fact distance at most
three from $K_n$ gives $-p\cdot Y\le3|Y|$. Multiplication by this function
is therefore continuous on Schwartz space, including as $Y$ approaches
a tied gap. This proves the partial boundary formula
\begin{equation}
 W_{n,Y}^\xi(f)=\widehat W_n^\xi
 \left(\chi e^{-p\cdot Y}(2\pi)^{-4n}
                         \int e^{ip\cdot X}f(X)\,dX\right).
 \label{ch16:partial-boundary}
\end{equation}
It has boundary $W_n^\xi$ at $Y=0$. The cutoff is essential here: the
exponential alone need not be a Schwartz multiplier off the spectral
cone. Tube Lorentz covariance follows from boundary covariance; changed
cutoffs agree near $K_n$ and hence give the same pairing.

\begin{proposition}[Adjacent spacelike exchange]
\label{ch16:locality}
Let $\sigma$ exchange slots $j$ and $j+1$. On the open set where
$X_{j+1}-X_j$ is spacelike,
\[
 W_n^\xi(\boldsymbol A;X)
       =W_n^\xi(\sigma\boldsymbol A;\sigma X)
\]
as distributions. This statement permits arbitrary exterior insertions.
\end{proposition}
\begin{proof}
Keep positive scalar imaginary gaps except for one tied pair
$y_j=y_{j+1}$. Write its spacelike real difference as $(v,t)$, and set
$e=v/|v|$. The imaginary part of its image under a boost of parameter
$i\alpha$ in direction $e$ is
$\sin\alpha\,(te,|v|)$, a strictly forward vector for small positive
$\alpha$. All other gaps stay strictly forward by continuity. Thus a
small positive imaginary boost enters the original tube, and a negative
one followed by exchange enters the exchanged tube.

Fix one sufficiently small nonzero boost parameter first. The transformed
tied-gap point is strictly inside the relevant tube, so a whole complex
neighborhood of the original point is mapped into that tube. Pulling
back the holomorphic function, with its entire finite-dimensional label
action, therefore defines a genuine germ on that neighborhood. Use the
opposite sign and the exchanged tube for the second germ. Real Lorentz covariance and the identity
theorem show that each germ extends its original tube function; the
whole short imaginary-boost path stays in that tube at neighboring
interior points.

At vanishing real times these germs coincide. Their restrictions to the
Euclidean slice $(it_k,x_k)$ equal the same smooth, permutation-symmetric
function on its two ordered sides. At the tied pair the spatial points
are distinct, so this Euclidean function is smooth there. Consequently
all derivatives of the two restrictions agree at the tie. The slice
is totally real and its tangent space spans the ambient complex vector
space. The holomorphic Taylor coefficients therefore agree, proving
identity of the germs.

For general real times multiply every real time by $s\in[0,1]$.
The distinguished difference remains spacelike along this path and the
other imaginary gaps remain strict. A finite chain of overlapping germ
neighborhoods propagates the equality from $s=0$ to $s=1$ by uniqueness.
Take the continuous partial boundaries \eqref{ch16:partial-boundary},
and then remove the other imaginary gaps. This proves the identity
against test functions compactly supported in the spacelike open set. Smooth
compact cutoffs converging in Schwartz space give the stated distributional
version. No permutation of coincident Lorentzian boundary values has
been assumed.
\end{proof}

\section{Field words and the reflection space}
\label{ch16:fields-section}

A letter is a pair $(A,f)$ with $A\in F$ and $f\in\mathscr S(\RR^4)$.
The star of a word reverses its order and replaces every letter by
$(A^*,\overline f)$. Let $\mathcal P$ be their algebraic span, including
the empty word $1$. The distributions $W_n^\xi$ define a functional
$w_\xi$ on $\mathcal P$ by evaluation on product test functions.

The essential positivity statement follows inside the phase reflection
space. Give a word strictly increasing positive imaginary heights,
smear its real variables, and use the slab product to obtain a vector
$h_y^\xi(b)$. Its Gram pairing with another such vector is a tested
continued correlation, with the reflected first word on the left.
Proposition~\ref{ch16:boundary} permits the two height lists to tend to
zero independently, so
\[
 \langle h_y^\xi(b),h_z^\xi(c)\rangle
                      \longrightarrow w_\xi(b^*c).
\]
Expanding $\|h_y^\xi(b)-h_z^\xi(b)\|^2$ proves that these vectors are
Cauchy. Their limits $v_\xi(b)$ have Gram form $w_\xi(b^*c)$, which is
therefore positive. The same common seminorm bounds extend positivity
to finite sequences of joint Schwartz test functions by tensor-product density.

The standard word reconstruction \cite{Wightman1956} now gives
\[
 \mathcal N_\xi=\{b:w_\xi(b^*b)=0\},\quad
 \mathcal D_\xi=\mathcal P/\mathcal N_\xi,\quad
 \mathcal H_\xi=\overline{\mathcal D_\xi},\quad\Omega_\xi=[1],
 \qquad \Phi_\xi(A,f)[b]=[(A,f)b].
\]
For clarity, its null-space check is worth recording. If $b$ is null,
Cauchy--Schwarz says it pairs to zero with every word. The concatenation
identity moves a letter to the other side of the pairing, so its product
with $b$ also pairs to zero with every word, including itself. Thus all
fields are well defined on the same invariant dense domain. The same
identity gives
\[
 \Phi_\xi(A^*,\overline f)\subset\Phi_\xi(A,f)^*.
\]
They are closable, and the joint distribution bound gives continuity of
$f\mapsto\Phi_\xi(A,f)v$ in Hilbert norm for every $v\in\mathcal D_\xi$.
We do not assert that these generally unbounded fields are bounded or
that their self-adjoint closures exist without further input.

\begin{theorem}[Fields and unitary equivalence with the reflection-positive space]
\label{ch16:fields}
In every phase the fields above are local, covariant under a strongly
continuous unitary representation of the proper orthochronous Poincar\'e
group, and have joint translation spectrum in $\overline V_+$.
The map
\[
 J_\xi^{\rm OS}[b]=v_\xi(b)
\]
extends to a unitary map from $\mathcal H_\xi$ onto the phase reflection
space. It sends vacuum to vacuum, intertwines spatial translations and
heat semigroups, and the field labeling is faithful.
\end{theorem}
\begin{proof}
Covariance preserves the word Gram form, hence defines unitaries on its
quotient. Schwartz continuity proves strong continuity first on words
and then on their completion. Locality follows by applying
Proposition~\ref{ch16:locality} inside arbitrary word matrix elements.
The joint spectral theorem for the commuting translation unitaries,
together with \eqref{ch16:cone-support} applied to such matrix elements,
places their spectrum in $\overline V_+$.

The Gram identity makes $J_\xi^{\rm OS}$ an isometry. Its range is closed,
but density must still be checked. For a smooth compactly supported
joint frequency test function $b$, multiplication by
$e^{-\sum_j y_jp_j^4}$ is harmless. Equation~\eqref{ch16:laplace} identifies
the norms and mixed pairing of the image of the resulting boundary word
with the smeared positive-height vector at heights $y$. They are the
same vector. Test functions with compact Fourier support are dense in Schwartz space;
therefore the closed range contains every smeared positive-height vector.
Approximate identities, followed by a perturbation of tied times and
permutation of separated Euclidean points, give all Euclidean word
vectors. These span the reflection space. The same positive-height
identities prove heat and spatial intertwining.

Finally a label acting as zero has every boundary correlation zero.
The Laplace identity, Euclidean symmetry and continuity imply that all
its separated Euclidean correlations vanish. The common null-space
identity of Chapter~\ref{ch15:phases}, and the injectivity of the matrix of correlation pairings used
there, force this label to be zero in $F$.
\end{proof}

\section[Vacuum uniqueness in each phase]{The low-energy estimate and uniqueness of the vacuum in each phase}
\label{ch16:vacuum-section}

The phase measure has full support, and the zero-energy space is a line
for almost every phase by \ref{ch15:ae-ground}. To remove the exceptional
set requires a bound uniform in $\xi$. We prove it from locality and the
spectral cone. This argument controls energies near zero, not a positive
mass gap.

For a finite sum $A$ of field words with compactly supported smooth test functions,
let $\mu_\xi$ be the joint spectral measure of $A\Omega_\xi$, and let
$\mu_\xi^*$ be that of $A^*\Omega_\xi$. Their total masses have a common
bound $M$ by Proposition~\ref{ch16:boundary}. Enclose all supports in a
fixed spacetime box. Locality gives some $R<\infty$, independent of phase,
such that
\begin{equation}
 a_\xi(y,t)=\langle\Omega_\xi,[A^*,\alpha_{(y,t)}(A)]\Omega_\xi\rangle=0
                  \quad (|y|>R,\ |t|\le1),
 \label{ch16:local-commutator}
\end{equation}
where $\alpha_z(A)=U_\xi(z)AU_\xi(z)^*$ on the word domain.
The spectral expression is the difference of the Fourier transforms of
$\mu_\xi$ and the reflected measure $\mu_\xi^*$.

Take the odd time average
\[
 \bar a_\xi(y)=\int_0^1
       \frac{a_\xi(y,t)-a_\xi(y,-t)}{2i}\,dt.
\]
It is bounded by $M$ and supported in $\overline B_R$. Its spatial
Fourier measure is the \emph{sum} of the pushforwards of
$\varpi(E)\,d\mu_\xi(p,E)$ and $\varpi(E)\,d\mu_\xi^*(p,E)$ under
$p$ and $-p$, respectively, where
\[
 \varpi(E)=\frac{1-\cos E}{E},\qquad\varpi(0)=0.
\]
Both measures are positive; on $0\le E\le1$,
$\varpi(E)\ge E/\pi$. Choose a smooth momentum bump $b_u$, equal to one
on $B_u$, supported in $B_{2u}$, and between zero and one. Its inverse
Fourier transform is bounded by $Cu^3$. On the shell $u/2<E\le u$,
cone support implies $|p|\le u$, so positivity and Fourier inversion yield
\[
 c u\,\mu_\xi\{u/2<E\le u\}
      \le \left|\int\bar a_\xi(y)\widehat b_u(y)\,dy\right|
      \le CMR^3u^3.
\]
Dyadic summation proves the first assertion below.

\begin{theorem}[Uniform low-energy bound]
\label{ch16:vacuum}
For every such $A$ there is $C_A$, independent of phase, such that
\begin{align}
 \mu_\xi\{0<E\le u\}&\le C_Au^2 &&(0<u\le1),\label{ch16:low-energy}\\
 0\le\langle A\Omega_\xi,(e^{-sH_\xi}-P_{\Omega,\xi})A\Omega_\xi\rangle
       &\le 2C_A' s^{-2} &&(s\ge1),\label{ch16:heat-decay}
\end{align}
where $P_{\Omega,\xi}$ is the projection onto $\CC\Omega_\xi$.
In every phase, $\ker H_\xi=\CC\Omega_\xi$; the same holds in its
reflection space.
\end{theorem}
\begin{proof}
Initially use the full zero-energy projection $Q_\xi$, without assuming
that its range is a line. Increase the constant in
\eqref{ch16:low-energy} so the estimate holds for all $u>0$ by the total
mass bound. Tonelli's theorem gives
\[
 \langle A\Omega_\xi,(e^{-sH_\xi}-Q_\xi)A\Omega_\xi\rangle
 =\int_{E>0}e^{-sE}\,d\mu_\xi
 \le C_A's\int_0^\infty e^{-su}u^2\,du=2C_A's^{-2}.
\]
Theorem~\ref{ch16:fields} transfers almost-everywhere vacuum uniqueness
from the reflection spaces, so $Q_\xi=P_{\Omega,\xi}$ almost everywhere.
For fixed $s>0$ the vacuum-subtracted heat pairing is continuous in $\xi$:
it is the joint distribution at a fixed weakly ordered height with gap
$s$ between its two word halves. This follows either by the independent
height limits or by \eqref{ch16:partial-boundary}. Thus its bound by
$2C_A's^{-2}$ holds everywhere, since any nonempty open violation set
would have positive $\nu$-measure.

Let $s\to\infty$. Strong heat convergence gives
$\|Q_\xi A\Omega_\xi\|^2=|\langle\Omega_\xi,A\Omega_\xi\rangle|^2$.
Since $Q_\xi$ fixes the unit vacuum, its value on each such word vector
lies in the vacuum line. Compactly smeared words are dense: multiply
Schwartz test functions by expanding smooth cutoffs and use joint Gram continuity.
It follows that $Q_\xi=P_{\Omega,\xi}$ on the full space. The onto heat
intertwiner gives the last assertion.
\end{proof}

\section{The weak commutant and local stress charges}
\label{ch16:charges-section}

Irreducibility should be stated in a form appropriate to unbounded fields.
Their common domain is $\mathcal D_\xi$, the span of field words applied
to the vacuum. A bounded operator $B$ belongs to their weak commutant if
\begin{equation}
 \langle\Phi^*u,Bv\rangle=\langle u,B\Phi v\rangle,
 \qquad u,v\in\mathcal D_\xi,
 \label{ch16:weak-commutant}
\end{equation}
for every smeared field $\Phi$. The adjoint inclusion proved above
ensures that these $u$ are in its adjoint domain; no preservation of
$\mathcal D_\xi$ by $B$ is assumed.

\begin{proposition}
\label{ch16:irreducible}
The weak commutant in every phase consists of scalar operators.
\end{proposition}
\begin{proof}
Iterate \eqref{ch16:weak-commutant} to move words through matrix elements.
For a word $A$ put $u=A\Omega_\xi$, $v=A^*\Omega_\xi$. The functions
\[
 p(z)=\langle u,e^{izH_\xi}B\Omega_\xi\rangle,\qquad
 q(z)=\langle B^*\Omega_\xi,e^{izH_\xi}v\rangle
\]
are bounded and holomorphic in the upper half-plane, continuous to its
boundary, and satisfy $p(t)=q(-t)$. Pasting across the real line and
using Morera's theorem gives a bounded entire function, hence a constant.
The heat limit and Theorem~\ref{ch16:vacuum} imply
$B\Omega_\xi=b\Omega_\xi$. Moving words again gives
$\langle A'\Omega_\xi,BA\Omega_\xi\rangle
=b\langle A'\Omega_\xi,A\Omega_\xi\rangle$.
Density proves $B=bI$. Scalars plainly satisfy the condition.
\end{proof}

Let $T_{ab}$ be the normalized Euclidean stress tensor of
Proposition~\ref{ch14:normalized-stress}. Define
\[
 \zeta_i=1\ (i\le3),\qquad\zeta_4=-i,
 \qquad T^{\rm M}_{ab}=-\zeta_a\zeta_bT_{ab}.
\]
We use the same notation for its smeared relativistic field. Thus
$T^{\rm M}_{44}=T_{44}$, $T^{\rm M}_{4i}=iT_{4i}$, and
$T^{\rm M}_{ij}=-T_{ij}$. The adjoint involution, covariance and the
continued derivative identity give a symmetric Hermitian Lorentz tensor
with
\[
 \sum_a\eta_{aa}T^{\rm M}_{ab}(\partial_a f)=0
                      \quad\hbox{on }\mathcal D_\xi.
\]
Indeed continuation of a Euclidean derivative label is
$\Phi_\xi(\partial_a A,f)=-\zeta_a^{-1}\Phi_\xi(A,\partial_a f)$.

For an affine Killing vector $K(x)=k+Vx$, where
$V^{\mathsf T}\eta+\eta V=0$, let $\Sigma_K$ be its infinitesimal
Lorentz label action. Choose real $f\in C_c^\infty(\RR)$ of integral
one and $c\in C_c^\infty(\RR^3)$, and define the local charge
\[
 Q_K(f,c)=\sum_b T^{\rm M}_{4b}\big(K^b(c\otimes f)\big).
\]
Suppose $g\in C_c^\infty(\RR^4)$ and $c=1$ on a ball such that every
point $(y,t)$ outside that ball with $t\in\operatorname{supp}f$ is
spacelike to $\operatorname{supp}g$. Then
\begin{equation}
 i[Q_K(f,c),\Phi_\xi(A,g)]
   =\Phi_\xi(\Sigma_K A,g)
          -\sum_a\Phi_\xi(A,\partial_a(K^ag))
                       \quad\hbox{on }\mathcal D_\xi.
 \label{ch16:charge-law}
\end{equation}
This is a local charge identity, not a claim that the unsmeared spatial
integral of stress exists as an operator.

To justify its continuation, first integrate the Euclidean shell identity
with a spatial cutoff $c$ on the two sides of the distinguished insertion.
Conservation moves the derivative to $c$; that term stays separated from
the insertion. The two current integrals have one-sided limits whose
jump is $\Sigma_K A+K\cdot\partial A$. Retain strictly ordered positive
heights for every spectator and put their complex times on a path
$z_j(\lambda)=iy_j+\lambda t_j$. Place the current immediately above
or below the distinguished time, at distance $\varepsilon$.

On $\operatorname{supp}\nabla c$ the spacelike holomorphic germs from
Proposition~\ref{ch16:locality} bound the current integrals uniformly
as $\varepsilon\downarrow0$, locally along this path. Conservation bounds
the difference between two such integrals by $C|\varepsilon-\varepsilon'|$.
Their limits are therefore holomorphic. At small imaginary $\lambda$
the jump is the Euclidean one; analytic continuation gives the same
jump at $\lambda=1$. Smearing with $f$ and $g$ and taking the joint
partial boundaries proves \eqref{ch16:charge-law}; integration by parts
produces $\partial_a(K^ag)$. General word spectators follow from the
common Schwartz bounds. Locality removes any change in $c$ outside the
stated ball. This retains precisely the normalization established by the
Euclidean shell flux.
